
\begin{definition}
\label{XI.4.4}
Sea~$G$ un $S$-preesquema. Se denota por~$\H^0(S,G)$ el conjunto de las
secciones de~$G$ sobre~$S$, que se considerará como grupo cuando~$G$ sea un
$S$-grupo. En tal caso, se denota por~$\H^1(S,G)$ el conjunto de las clases,
salvo isomorfismo, de fibrados homogéneos principales sobre~$S$ de grupo
\ifthenelse{\boolean{orig}}
{$S$}
{$G$},
considerando~$\H^1(S,G)$ dotado del «punto marcado» que corresponde a los
fibrados triviales\footnote{\label{note-296}Esta notación solo es coherente
\ifthenelse{\boolean{orig}}
{con respecto a}
{con respecto a}
las notaciones cohomológicas generales (SGA~4~V) cuando se dispone de
criterios de efectividad del descenso, que apenas están garantizados salvo si
$G$ es afín (o tan solo cuasiafín, \cf (\Ref{VIII.7.9})).}.
\end{definition}

Así, $\H^0(S,G)$ es un funtor de la variable $S$-preesquema~$G$, con valores
en la categoría de los conjuntos. Este funtor es exacto por la izquierda y,
con mayor razón, conmuta con los productos finitos, lo cual implica, en
efecto, que transforma grupos en grupos y grupos conmutativos en grupos
conmutativos. De manera análoga, $\H^1(S,G)$ es un funtor de la variable
$S$-grupo \emph{afín}~$G$, con valores en la categoría de los conjuntos por
medio de la formación de fibrados asociados; se comprueba fácilmente que
este funtor conmuta con los productos finitos. En particular, transforma los
grupos de la categoría de los $S$-grupos afines, \ie los $S$-grupos
\emph{afines conmutativos}, en grupos de la segunda categoría, e incluso en
grupos conmutativos (pues los grupos de la primera categoría son
conmutativos). Así, \emph{si~$G$ es un $S$-grupo afín conmutativo,
$\H^1(S,G)$ es un grupo conmutativo}, y un homomorfismo~$G\to H$ de
$S$-grupos afines conmutativos da lugar a un homomorfismo de grupos
$\H^1(S,G)\to\H^1(S,H)$.

Para simplificar, en lo sucesivo nos limitaremos a considerar $S$-grupos
\emph{afines y conmutativos}. Sea
$$
0\to G'\lto{u} G\lto{v}G''\to 0
$$
una sucesión de morfismos de tales grupos; \emph{diremos que esta sucesión es
exacta si~$vu=0$} (lo que permite considerar~$G$ como un preesquema
sobre~$G''$, con grupo de operadores por la derecha~$G'_{G''}$) \emph{y
si~$G$ es un fibrado homogéneo principal
\marginpar{297}%
sobre~$G''$ de grupo~$G'_{G''}=G'\times_S G''$}. Esto implica, en particular,
que~$u\colon G'\to G$ es un núcleo de~$v$ y, con mayor razón, implica la
exactitud de la sucesión
$$
\ifthenelse{\boolean{orig}}
{0\to\H^0(X,G')\to\H^0(X,G)\to\H^0(X,G'').}
{0\to\H^0(S,G')\to\H^0(S,G)\to\H^0(S,G'').}
$$
Esto implica, además, la posibilidad de definir una aplicación
$$
\ifthenelse{\boolean{orig}}
{\partial\colon\H^0(X,G'')\to\H^1(X,G')\quoi,}
{\partial\colon\H^0(S,G'')\to\H^1(S,G')\quoi,}
$$
asociando a toda sección de~$G''$ sobre~$S$, \ie a todo $S$-morfismo
$f\colon S\to G''$, el fibrado homogéneo principal~$P_f$ de grupo
$G'\simeq f^*(G'_{G''})$ sobre~$S$, imagen inversa del fibrado homogéneo
principal~$G$ sobre~$G''$. Desde el punto de vista de los $S$-preesquemas, no
es, pues, sino la imagen inversa por~$v\colon G\to G''$ del subpreesquema
imagen de~$S$ por la inmersión~$f$, y las acciones de~$G'$ sobre~$P_f$ son
inducidas por las acciones por la derecha de~$G'$
\ifthenelse{\boolean{orig}}
{sobre~$G$).}
{sobre~$G$.}
Dejamos asimismo al lector la comprobación de la proposición siguiente, que
no presenta más dificultad que la de su redacción:

\begin{proposition}
\label{XI.4.5}
La aplicación
\ifthenelse{\boolean{orig}}
{$\partial\colon \H^0(X,G'')\to\H^1(X,G')$}
{$\partial\colon \H^0(S,G'')\to\H^1(S,G')$}
es un homomorfismo de grupos. La sucesión de homomorfismos siguiente es
exacta:
\begin{multline*}
\ifthenelse{\boolean{orig}}
{0\to\H^0(X,G')\lto{u^0}\H^0(X,G)\lto{v^0}\H^0(X,G')\lto{\partial}\H^1(X,G'')\\
\lto{u^1}\H^1(X,G) \lto{v^1}\H^1(X,G'')}
{0\to\H^0(S,G')\lto{u^0}\H^0(S,G)\lto{v^0}\H^0(S,G'')\lto{\partial}\H^1(S,G')\\
\lto{u^1}\H^1(S,G) \lto{v^1}\H^1(S,G'')}
\end{multline*}
(donde los homomorfismos distintos de~$\partial$ proceden de la ley
funtorial de~$\H^0$ \resp de~$\H^1$).
\end{proposition}

\begin{remarques}
\label{XI.4.6}
El punto de vista expuesto aquí para el estudio de los fibrados homogéneos
principales está visiblemente inspirado en Serre~\cite{XI.7}, cuya consulta
será de gran provecho para el lector. Cuando se desea un formalismo que se
aplique también a $S$-grupos estructurales cuasiproyectivos sobre~$S$ (de
manera que abarque, en particular, los esquemas abelianos proyectivos),
conviene modificar~\Ref{XI.4.1} exigiendo que~$S'$ sea suma de
preesquemas~$S'_i$ finitos y localmente libres sobre abiertos~$S_i$ de~$S$
que recubran~$S$. Los desarrollos precedentes siguen siendo entonces
válidos, incluido en particular~\Ref{XI.4.5}, sustituyendo en todas partes la
hipótesis de afinidad por la de cuasiproyectividad e interpretando de manera
correspondiente la definición dada
\ifthenelse{\boolean{orig}}
{anteriormente}
{anteriormente}
de una sucesión exacta de $S$-grupos. En efecto, basta sustituir la referencia
a~(\Ref{VIII.2.1})
\marginpar{298}
por~(\Ref{VIII.7.7}): los
\ifthenelse{\boolean{orig}}
{morfismos utilizados}
{morfismos utilizados}
$S'\to S$ son morfismos de descenso efectivo para la categoría fibrada de
los preesquemas cuasiproyectivos sobre otros preesquemas. Sin embargo, debe
advertirse que esta segunda noción de fibrado homogéneo principal es más
restrictiva que la primera de~\Ref{XI.4.1}.
\end{remarques}

\sganum
\label{XI.4.7}
Se obtiene una noción aún más restrictiva de fibrado homogéneo principal
exigiendo que~$S$ esté recubierto por abiertos~$S_i$ tales que, para todo~$i$,
$P|S_i$ sea un fibrado con operadores trivial bajo~$G|S_i$; se dirá entonces
que~$P$ es un fibrado homogéneo principal \emph{localmente trivial}.
\index{homogéneo principal localmente trivial (fibrado)|hyperpage}%
Las clases de estos fibrados, para~$G$ dado, forman una parte de
\ifthenelse{\boolean{orig}}
{$\H^1(X,G)$}
{$\H^1(S,G)$},
que está en correspondencia biyectiva con
\ifthenelse{\boolean{orig}}
{$\H^1(X,\cal{O}_X(G))$}
{$\H^1(S,\cal{O}_S(G))$},
donde
\ifthenelse{\boolean{orig}}
{$\cal{O}_X(G)$}
{$\cal{O}_S(G)$}
es el haz (en el sentido ordinario) de las secciones
de~$G$ sobre~$S$, \cf \cite{XI.2}. Para que estos~$\H^1$ sigan dando lugar a
una sucesión exacta de cohomología~\Ref{XI.4.5}, es preciso suponer,
evidentemente, que la sucesión~$0\to G'\to G\to G''\to 0$ sea exacta en el
sentido razonable para este nuevo contexto, \ie que~$G$ sea un fibrado
localmente trivial sobre~$G''$, de grupo~$G'_{G''}$; esto significa también
que~$u\colon G'\to G$ es un núcleo de~$v\colon G\to G''$ y que~$G$ admite
localmente una sección sobre~$G''$.

\sganum
\label{XI.4.8}
Es, evidentemente, muy deseable prolongar la sucesión exacta~\Ref{XI.4.5}
introduciendo los grupos de cohomología superiores
\ifthenelse{\boolean{orig}}
{$\H^i(X,G)$}
{$\H^i(S,G)$}.
Esto es posible situándose en el marco de la «cohomología de Weil»: se
considera la categoría~$\cal{B}$ de los preesquemas cuasicompactos sobre~$S$,
\ifthenelse{\boolean{orig}}
{dotada}
{dotada}
del conjunto~$\cal{M}$ de los morfismos fielmente planos y cuasicompactos,
que se llamarán morfismos localizantes. Un «haz de Weil» abeliano sobre~$S$
(o, mejor, sobre~$(\cal{B},\cal{M})$) es entonces un funtor contravariante
$\cal{F}$ de~$\cal{B}$ en la categoría de los grupos abelianos, que transforma
sumas en productos y una sucesión
\begingroup
\shorthandoff{<>}% Spanish babel trata de otro modo los delimitadores de curvatura de Xy-pic.
${\def\labelstyle{\scriptstyle}\xymatrix@C=.5cm{T''=T'\times_TT'\ar@<+2pt>[rr]^-{\pr_1,\pr_2}\ar@<-2pt>[rr] &&T'\ar[r]^f & T}}$,
con~$f\in\cal{M}$, en un diagrama \emph{exacto} de conjuntos
$\xymatrix@C=.5cm{\cal{F}(T)\ar[r] & \cal{F}(T')
\ar@<2pt>[r]\ar@<-2pt>[r] & \cal{F}(T'')}$.%
\endgroup\ 
Los haces de Weil forman una categoría abeliana con límites inductivos
exactos que admite un generador y, por tanto, que posee suficientes objetos
inyectivos~\cite{XI.1}. Los funtores derivados por la derecha del funtor
$\Gamma(\cal{F})=\cal{F}(S)$ se denotan entonces por~$\H^i(S,\cal{F})$. Por
otra parte, todo $S$-grupo conmutativo define evidentemente un haz de Weil
(\Ref{VIII.5.2}), cuyos $\H^0$ y~$\H^1$ no son sino~$\H^0(S,G)$ y
$\H^1(S,G)$, lo que permite definir de manera razonable los demás
$\H^i(S,G)$. Se
\marginpar{299}
demuestra, además, que una sucesión exacta de $S$-grupos define una sucesión
exacta de haces de Weil, lo que permite recuperar y prolongar la sucesión
exacta~\Ref{XI.4.5}\footnote{Para un estudio sistemático de este punto de
vista, \cf SGA~4~I a~IX.}.

\sganum
\label{XI.4.9}
Sería conveniente desarrollar las variantes no conmutativas de~\Ref{XI.4.5},
como en~\cite{XI.2}. Para un desarrollo sistemático, en el marco adecuado,
de las diversas nociones cohomológicas esbozadas en el presente número,
remitimos a un trabajo entonces en preparación de
J\ptbl\ \textsc{Giraud}\footnote{\Cf J\ptbl\ \textsc{Giraud},
\ifthenelse{\boolean{orig}}
{\emph{Álgebra homológica no abeliana}, de próxima aparición en}
{\emph{Cohomología no abeliana},}
Springer-Verlag 1971.}.
